\chapter{Model and Simulation Method}

\section{Model construction}
Describe geometry, material properties, excitations, ports, boundary conditions,
monitors, mesh controls, solver type, stopping criteria, and post-processing. A second
engineer should be able to reconstruct the model from the information presented here
and in the simulation record.

\section{Parameter control}
Maintain a single source of truth for design parameters. Table~\ref{tab:parameters}
illustrates a compact public parameter record.

\begin{table}[H]
\centering
\caption{Example public parameter record.}
\label{tab:parameters}
\begin{tabular}{llll}
\toprule
Parameter & Symbol & Example value & Evidence source \\
\midrule
Reference frequency & $f_0$ & \SI{2.50}{\giga\hertz} & Requirement \\
Normalized wavelength & $\lambda_0$ & \SI{119.9}{\milli\meter} & Analytical script \\
Nominal spacing & $d$ & $0.5\lambda_0$ & Analytical design \\
Relative phase & $\alpha$ & \ang{90} & Steering case \\
\bottomrule
\end{tabular}
\end{table}

\section{Run manifest}
Record each meaningful numerical run with a stable identifier, model revision, solver,
mesh or discretization setting, changed parameter, completion status, and exported
artifacts. The starter file \texttt{simulation\_records/run\_manifest.csv} provides a
machine-readable structure.

\begin{simulatedevidencebox}
A screenshot alone is not a complete simulation record. Pair every published plot with
a run identifier and enough configuration information to reconstruct its origin.
\end{simulatedevidencebox}
